\section{Discussion}
\label{sec:discussion}

\paragraph{What any agent framework should adopt.} The results suggest a short checklist that is independent of \method{}'s code. (1) Establish controls once, at entry, as an execution scope that every boundary---threads, generators, resume, fan-out, exports---re-enters, and compose nested scopes by a meet. (2) Never persist live handles; re-derive controls on resume. (3) Decide every effect at one deny-first admission point, and evaluate rate limits last so denials do not consume capacity. (4) Key limits by principal, not by agent or thread, so that fan-out cannot multiply a quota, and shard by that key. (5) Make the fail mode local configuration (Proposition~\ref{prop:failmode}). (6) Emit exactly one audit record per decision, including denials and outages. (7) Treat unmeasured as distinct from passed, and report governance scores together with ablations or trivial-runner baselines; a score that a deny-everything runner nearly matches says little by itself.

\paragraph{Why governance, and not task accuracy, is where the architecture shows.} Task benchmarks such as WorkBench are dominated by the model: published full-set runs range from 406 of 690 for the model we used to 674 of 690 for the strongest, a spread that dwarfs the differences among our arms. A governed platform's value is that controls hold when the model, the topology or the deployment changes, which is exactly what task accuracy does not measure. Conversely, governance benchmarks are cheap to run and deterministic, which makes them easy to overfit; mechanism attribution and discriminating scenarios are the counterweight.

\subsection{Threats to validity}
\label{sec:threats}
\paragraph{Benchmark-driven development.} The control plane was written during this study, after reading AgentGovBench's scenarios. A system built to a test suite can pass it without generalizing. We mitigate this in three ways, none of which removes the threat: every allow/deny decision is made by library code that has its own unit and property tests independent of the benchmark; ablations show that each passing category depends on a mechanism rather than on scenario-specific logic; and eight supplemental scenarios probe behavior the official suite does not. The supplemental scenarios are ours, so an independent suite would be stronger evidence.
\paragraph{Harness choices.} The virtual clock and threaded fan-out differ from the upstream runners, and both are disclosed. Threaded fan-out is stricter. The virtual clock matters only for the one recovery scenario, where the upstream reference runner sleeps through the same outage.
\paragraph{Comparators.} Published AgentGovBench rows were produced by their authors against a live gateway, not rerun by us; their declared-unsupported scenarios are counted as published. The trivial runners are ours.
\paragraph{Scope of the non-governance benchmarks.} The long-context benchmark is synthetic and measures assembly, not use by a model; the injection data set is small (\InjTestN{} test prompts) and partly non-English; HaluEval-QA has the length artifact we document; the scaling results are from one machine; and the WorkBench study uses one model, one repetition and subsets of 12 and 60 tasks.

\subsection{Limitations}
We did not measure the utility cost of attenuation, which MasDrift~\cite{xu2026masdrift} and Bounded Agents~\cite{muruaga2026boundedagents} show can be substantial; AgentGovBench has no model and therefore cannot. The audit log is not tamper-evident, approvals are not bound to action digests or consumed once, and there is no information-flow tracking; the Agent Governance Toolkit, Bounded Agents and CaMeL/FIDES are ahead on each. Process-local scopes are not an isolation boundary against malicious Python tools; the platform's in-process sandbox is for trusted code only. Cooperative cancellation does not stop a tool already running. Distributed budgets, fenced ownership of shared state, durable reconciliation of external effects (no exactly-once guarantee), authenticated approval resume for exported tools, and opaque CrewAI internals remain open, as does multi-agent planning quality, which we did not evaluate. AgentDojo, $\tau^2$-bench and MultiAgentBench require model-driven runs that we did not perform. The trained injection gate is a lexical model and should be replaced by or combined with a model-based detector where available.

\paragraph{Broader impact.} Explicit, testable controls can make it safer for small teams to deploy agents that take actions. The same results could be over-read: a perfect governance score does not certify a deployment, and filters do not make untrusted content safe. We report the conditions under which each result holds and the mechanisms it depends on, so that readers can judge transfer to their setting.

\section{Conclusion}
Controls attached to a request must survive every boundary a platform creates, and the way to make them survive is architectural: represent them as a composable execution scope and decide every effect at one deny-first admission point. We formalized this as the Governed Execution Contract, proved its central properties, implemented it in \method{}, and showed a perfect, repeatable AgentGovBench score whose passing scenarios are attributed to specific mechanisms. The same analysis shows that governance benchmarks need ablations, trivial baselines and discriminating scenarios before their scores can be read as evidence of enforcement, and we supply those. Across the other layers the platform is competitive or better than common defaults, and where it was not---injection filtering and paraphrase retrieval---the benchmarks told us what to fix.

\begin{ack}
AI assistance (Anthropic Claude) was used for implementation, literature search and citation checking, experiment orchestration and drafting. All experiments were executed by scripts in the repository, and every reported number is regenerated from raw evidence by \code{build\_evidence.py}. The author is responsible for the content. No external funding was received.
\end{ack}
